\section{Discussion}

The averaging denominator is part of the training objective. Global-token averaging weights domains by generated token counts; domain-token averaging fixes domain weights; domain-response averaging additionally gives responses equal weight within each domain. The comparison shows a capability tradeoff: stronger mathematics under global-token averaging accompanies weaker instruction following than under domain-response averaging.

Parameter concentration and vocabulary coverage describe different aspects of distillation. Both joint losses produce concentrated cumulative changes in BF16 model weights, while FP32 optimizer weights move more broadly. At a common state, the four local supervision choices yield very similar update concentration. Their online capability trajectories differ by domain and checkpoint. These observations motivate evaluating a distillation loss through both its update measurements and the capabilities it transfers.

The experiments cover one model family and one training seed. The local comparison uses four prefix banks with high top-$k$ probability coverage; the bootstrap intervals quantify evaluation-question uncertainty at fixed checkpoints.
